% ===== Block B12: 原文 p.101–115 | §5.3 边界附近的微分方程 / §5.4 单点椭圆曲线的共形块 / §5.5 Verlinde 公式 | 译手 tr-ch5a =====

\subsection{边界附近的微分方程}

% [原文 p.101]
本节我们要证明在 (4.26) 中构造的形式解收敛。

设 $\mathfrak{F}^{(0)}=(\pi:\cC\to B;\,s_1,\dots,s_N)$ 是 $N$-点稳定曲线的通用族，而 $\mathfrak{F}=(\pi:\cC\to B;\,s_1,\dots,s_N;\eta_1,\dots,\eta_N)$ 是在 $\mathfrak{F}^{(0)}$ 上添加形式邻域所得到的族。我们采用与 4.3 节相同的记号。设 $D$ 是族 $\mathfrak{F}^{(0)}$ 的判别轨迹，并设 $D=\bigcup_{i=1}^k D_i$ 是它分解为不可约分支的分解。本节中为简单起见假设 $k=1$。令 $E=D_1$。记 $\pi_E:\cC_E\to E$ 为 $\cC$ 到 $E$ 上的限制。设 $\tilde\pi_E:\tilde{\cC}_E\to E$ 由沿 $E$ 上二重点的轨迹对 $\pi_E:\cC_E\to E$ 作同时正规化而得到，而 $\sigma',\sigma'':E\to\tilde{\cC}_E$ 是对应于该轨迹 $E$ 上正规化二重点的截面。

% [Tutor Note: $k=1$ 的假设是可用的关键：判别轨迹 $D$ 只有一个不可约分支 $E$ 时，
%  纤维 $\cC_E$ 在 $E$ 的每个点处只有通常二重点（两个分支），于是引理 5.15 的
%  局部形式恰是"两个分支各配一个局部坐标 $x,y$"；若 $k\ge 2$，$D_i$ 之间会相交，
%  纤维上出现三重及以上的奇点，同时正规化不再是两个分支的并，$\alpha$ 要换成多重留数和。]

% [原文 p.102]
\[
\begin{array}{ccccc}
\tilde{\cC}_E & \longrightarrow & \cC_E & \longrightarrow & \cC^{(0)}\\[2pt]
& \swarrow{\scriptstyle \tilde\pi_E,\ s_j,\sigma',\sigma''}
  & \uparrow{\scriptstyle \pi_E,\ s_j,\sigma}
  & & \big\downarrow\\[2pt]
& E & & \longrightarrow & \cB
\end{array}
\]

% [Tutor Note: 原书此处是一张无编号的交换图表（不含在图的编号系统里，故不用 \figomit）。
%  这里按行用 array 重排，箭头读作：上面一行是 $\tilde\cC_E\to\cC_E\to\cC^{(0)}$ 的
%  自然映射链；$\tilde\pi_E$ 与截面 $s_j,\sigma',\sigma''$ 都是 $\tilde\cC_E$ 与 $E$ 之间的
%  映射（一上一下）；$\pi_E$ 与截面 $s_j,\sigma$ 是 $\cC_E$ 与 $E$ 之间的映射（一上一下）；
%  末列是 $\cC^{(0)}\to\cB$ 与 $E\to\cB$ 的投影。核对方法：$E$ 上每个二重点处
%  $\sigma',\sigma''$ 是正规化曲线上的两点，它们在 $\cC_E$ 中粘成一点 $\sigma(u)$。]

如命题 4.13 所证明的，族
\[
\tilde{\mathfrak{F}}^{(0)}_E:=\bigl(\tilde\pi_E:\tilde{\cC}_E\to E;\,s_1,\dots,s_N,\sigma',\sigma''\bigr)
\]
是 $(N+2)$-点（不一定连通）稳定曲线的局部通用族。

\begin{lemma}[5.15]
若 $B$ 取得充分小，则存在 $\sigma'(E)$（相应地 $\sigma''(E)$）在 $\tilde{\cC}_E$ 中的邻域 $X$（相应地 $Y$）的局部坐标 $(u_1,\dots,u_M,z)$（相应地 $(u_1,\dots,u_M,w)$），以及相对向量场
\[
\tilde\ell\in H^0\Bigl(\tilde{\cC}_E,\ \Theta_{\tilde{\cC}_E/E}\Bigl(*\sum_{j=1}^N s_j(E)\Bigr)\Bigr),
\]
它们满足以下条件。
\begin{enumerate}[label=(\arabic*),leftmargin=2.6em]
\item 截面 $\sigma'$ 与 $\sigma''$ 由如下映射给出：
\[
\begin{array}{rclll}
\sigma' & : & (u_1,\dots,u_M) & \mapsto (u_1,\dots,u_M,0)=(u_1,\dots,u_M,z),\\[2pt]
\sigma'' & : & (u_1,\dots,u_M) & \mapsto (u_1,\dots,u_M,0)=(u_1,\dots,u_M,w).
\end{array}
\]
\item $\tilde\ell\big|_X=\dfrac12\,z\dfrac{\partial}{\partial z}$,\qquad
      $\tilde\ell\big|_Y=\dfrac12\,w\dfrac{\partial}{\partial w}$.
\end{enumerate}
\end{lemma}

% [Tutor Note: 条件 (2) 里的系数 $1/2$ 不是任意的。它是把 $\tilde\ell$ 沿 $\sigma'(E)$
%  的线性项归一化：$\tilde\ell$ 在 $X$ 上写成 $a(u,x)\partial/\partial x$，条件 (2) 等价于
%  $a(u,x)=\frac12 x+O(x^2)$，这正是下文 $z$ 作为坐标的归一化 $\partial z/\partial x|_{x=0}=1$
%  所要求的；若系数取 $c\ne 1/2$，则 $z=x^{2c}$，$c\ne1/2$ 时 $z$ 不是 $x=0$ 处的坐标。]

\noindent\textit{证明.}\quad 设 $\nu:\tilde{\cC}_E\to\cC_E$ 是沿 $E$ 上二重点的同时正规化。设 $(u_1,\dots,u_M,x)$（相应地 $(u_1,\dots,u_M,y)$）是 $X$（相应地 $Y$）上满足上述引理中条件 (1) 的局部坐标。由于 $\nu$ 在 $\tilde{\cC}_E\setminus(\sigma'(E)\cup\sigma''(E))=\cC_E\setminus\sigma(E)$ 上是同构（恒等映射），由引理 1.33 的证明、特别是由 (1.60)，我们得到如下正合列
\[
0\to\Theta_{\cC_E/E}\to\nu_*\bigl(\Theta_{\tilde{\cC}_E/E}(-\sigma'(E)-\sigma''(E))\bigr)\xrightarrow{\ \alpha\ }\cO_E\to 0,
\]
其中 $\cO_E$-模同态 $\alpha$ 由下式给出：
\[
\Bigl(a(u,x)\frac{\partial}{\partial x},\,b(u,y)\frac{\partial}{\partial y}\Bigr)
\longmapsto \frac{\partial a(u,0)}{\partial x}+\frac{\partial b(u,0)}{\partial y}.
\]
注意 $\nu_*\bigl(\Theta_{\tilde{\cC}_E/E}(-\sigma'(E)-\sigma''(E))\bigr)$ 在点 $\sigma(u)$（$u\in E$）处的茎由一对局部全纯向量场 $\bigl(a(u,x)\frac{\partial}{\partial x},\allowbreak b(u,y)\frac{\partial}{\partial y}\bigr)$ 组成，其中 $a(t,0)=0$、$b(y,0)=0$，而 $\alpha$ 的定义与局部坐标的选取无关。该正合列诱导出正合列
\[
\begin{array}{rcl}
0 & \longrightarrow & H^0(\cC_E,\Theta_{\cC_E/E}(kS))\\[2pt]
  & \longrightarrow & H^0\bigl(\cC_E,\nu_*(\Theta_{\tilde{\cC}_E/E}(\nu^*(kS)-\sigma'(E)-\sigma''(E)))\bigr)\\[2pt]
  & \xrightarrow{\ \alpha\ } & H^0(E,\cO_E)\longrightarrow H^1(\cC_E,\Theta_{\cC_E/E}(kS))
\end{array}
\]

% [Tutor Note: 关于记号 $*$：$*\sum_j s_j(E)$ 表示允许沿各标记点取任意阶极点
%  （形式 Laurent 展开），因此 $H^0$ 是无穷维的；下面要靠 $k$ 足够大消去 $H^1$ 才能把
%  $\alpha(\tilde\ell)\equiv1$ 的局部解拼成整体解。这是本节的核心思路。]

% [原文 p.103]
对每个整数 $k$，其中我们置
\[
S=\sum_{j=1}^N s_j(E).
\]
当 $k$ 充分大时，我们有
\[
H^1(\cC_E,\Theta_{\cC_E/E}(kS))=0.
\]
于是由上述正合列，存在相对向量场
\[
\begin{array}{rcl}
\tilde\ell & \in & H^0\bigl(\tilde{\cC}_E,\Theta_{\tilde{\cC}_E/E}(\nu^*(kS)-\sigma'(E)-\sigma''(E))\bigr)\\[2pt]
& = & H^0\bigl(\cC_E,\nu_*(\Theta_{\tilde{\cC}_E/E}(\nu^*(kS)-\sigma'(E)-\sigma''(E)))\bigr),
\end{array}
\]
使得
\[
\alpha(\tilde\ell)\equiv 1.
\]
在上述给定的局部坐标下，$\tilde\ell$ 具有形式
\[
\tilde\ell=a(u,x)\frac{\partial}{\partial x}\quad\text{在 }X\text{ 上},\qquad
\tilde\ell=b(u,y)\frac{\partial}{\partial y}\quad\text{在 }Y\text{ 上},
\]
并且
\[
\frac{\partial a(u,0)}{\partial x}+\frac{\partial b(u,0)}{\partial y}\equiv 1.
\]
必要时加上来自 $H^0(\cC_E,\Theta_{\cC_E/E}(*S))$ 的元素，并把 $B$、$X$、$Y$ 取得更小，我们可以假设 $X$ 与 $Y$ 具有 Taylor 展开
\[
a(u,x)=\frac12 x+\sum_{n=2}^{\infty}a_n(u)x^n,\qquad
b(u,y)=\frac12 y+\sum_{n=2}^{\infty}b_n(u)x^n.
\]
现在由下式定义 $z=z(u,x)$ 与 $w=w(u,y)$：
\[
a(u,x)\frac{\partial z}{\partial x}=\frac12 z,\quad z(u,0)=0;\qquad
b(u,y)\frac{\partial w}{\partial y}=\frac12 w,\quad w(u,0)=0.
\]
于是，通过把 $X$ 与 $Y$ 取得更小，$(u_1,\dots,u_M,z)$ 与 $(u_1,\dots,u_M,w)$ 满足上述条件 (1) 与 (2)。\hfill$\square$

% [Tutor Note: 原文 $b(u,y)=\frac12 y+\sum_{n\ge2}b_n(u)x^n$ 中的 $x^n$ 应为 $y^n$
%  （$b$ 是 $y$ 的幂级数，与上一行 $a$ 中的 $x^n$ 对举可判；否则 $b$ 关于 $y$ 的展式不完整）。]

% [Tutor Note: 证明中 "a(t,0)=0, b(y,0)=0" 的变量混用是原文笔误：按 $\alpha$ 的定义
%  （对 $x$ 与 $y$ 分别取 $0$）应为 $a(u,0)=0$、$b(u,0)=0$；另外 $a(u,0)=0$ 与
%  $\tilde\ell$ 沿 $\sigma'(E)$ 属于 $-\sigma'(E)$ 极点除子一致。]

设 $\tilde{\mathfrak{F}}_E=(\tilde\pi_E:\tilde{\cC}_E\to E;\,s_1,s_2,\dots,s_N,\sigma',\sigma'';\eta_1,\eta_2,\dots,\eta_N,z,w)$ 是由 $\tilde{\mathfrak{F}}^{(0)}_E$ 得到的、带有形式坐标的 $(N+2)$-点亏格 $g$ 稳定曲线族。

设 $\ell_j(\xi_j)\dfrac{d}{d\xi_j}$ 是 $\tilde\ell$ 关于形式坐标 $\xi_j$ 的形式 Laurent 展开。于是 $\ell_j(\xi_j)\in\cO_E((\xi_j))$。置
\begin{equation}
\vec\ell=\Bigl(\ell_1(\xi_1)\frac{d}{d\xi_1},\dots,\ell_N(\xi_N)\frac{d}{d\xi_N}\Bigr). \tag{5.23}
\end{equation}
% [原文 p.104]
\figomit{5.1}{模空间边界附近的黎曼曲面。}

下面我们从族 $\tilde{\mathfrak{F}}^{(0)}_E$ 重新构造族 $\mathfrak{F}^{(0)}$。使用引理 5.15 的记号，我们可以假设
\[
\begin{array}{rcl}
X & = & \{\,P\in\tilde{\cC}_E\mid |z(P)|<1\,\}\\[2pt]
Y & = & \{\,P\in\tilde{\cC}_E\mid |w(P)|<1\,\}.
\end{array}
\]
对正数 $\eps<1$，置
\[
\begin{array}{rcl}
X_\eps & = & \{\,P\in\tilde{\cC}_E\mid |z(P)|<\eps\,\}\\[2pt]
Y_\eps & = & \{\,P\in\tilde{\cC}_E\mid |w(P)|<\eps\,\}.
\end{array}
\]
固定正数 $\eps_1<\eps_2<1$，并设 $\{U_\alpha\}_{3\le\alpha\le A}$ 是 $\tilde{\cC}_E\setminus(X_{\eps_2}\cup Y_{\eps_2})$ 的一个有限开覆盖，使得对任意 $\alpha=3,\dots,A$ 有
\[
U_\alpha\cap X_{\eps_1}=\emptyset,\qquad U_\alpha\cap Y_{\eps_1}=\emptyset.
\]
现在置
\[
\begin{array}{rcl}
D & = & \{\,\tau\in\C\mid |\tau|<1\,\}\\[2pt]
S_0 & = & \{\,(x,y,\tau)\in\C^3\mid xy=\tau,\ |x|<1,\ |y|<1,\ |\tau|<1\,\}\\[2pt]
S & = & S_0\times E\\[2pt]
Z & = & \{\,(P,\tau)\in\tilde{\cC}_E\times D\mid P\in\tilde{\cC}_E\setminus(X\cup Y)\\
  &   & \qquad\text{或 } P\in X \text{ 且 } |z(P)|>|\tau|\,\}\\[2pt]
W & = & \{\,(P,\tau)\in\tilde{\cC}_E\times D\mid P\in\tilde{\cC}_E\setminus(X\cup Y)\\
  &   & \qquad\text{或 } P\in Y \text{ 且 } |w(P)|>|\tau|\,\}.
\end{array}
\]
在 $Z\sqcup S\sqcup W$ 上如下引入等价关系 $\sim$。
\begin{enumerate}[label=(\arabic*),leftmargin=2.6em]
\item $Z\cap(X\times D)$ 中的点 $(P,\tau)$ 与 $S$ 中的点 $(x,y,\tau',u)$ 等价，当且仅当
\[
(x,y,\tau',u)=\Bigl(z(P),\ \frac{\tau}{z(P)},\ \tau,\ \tilde\pi_E(P)\Bigr).
\]
\item $W\cap(Y\times D)$ 中的点 $(P,\tau)$ 与 $S$ 中的点 $(x,y,\tau',u)$ 等价，当且仅当
\[
(x,y,\tau',u)=\Bigl(\frac{\tau}{w(P)},\ w(P),\ \tau,\ \tilde\pi_E(P)\Bigr).
\]
\item $Z$ 中的点 $(P,\tau)$ 与 $W$ 中的点 $(Q,\tau')$ 等价，当且仅当
\[
(P,\tau)=(Q,\tau').
\]
\end{enumerate}
现在置 $\cC=Z\sqcup S\sqcup W/\!\sim$。于是容易证明 $\cC$ 是复流形，并且存在自然的全纯映射 $\pi:\cC\to E\times D$。此外，由于我们

% [Tutor Note: 构造是标准的"管道装置（plumbing）"：$S_0=\{|x|,|y|<1,\ xy=\tau\}$ 是
%  光滑化后的结点邻域，$Z$、$W$ 分别是"只保留 $z$ 侧"与"只保留 $w$ 侧"的部分，
%  粘合条件 (1)(2) 让 $\tau$ 小的纤维把两块沿环面 $|z|=|\tau|^{1/2}$ 拼起来。
%  核对：$Z$ 中要求 $|z(P)|>|\tau|$、$W$ 中要求 $|w(P)|>|\tau|$，于是
%  $|z|=|\tau|$ 与 $|w|=|\tau|$ 的环带正好被剪掉换成了 $S$。]

% [原文 p.105]
可以假设 $s_j(E)$ 含于 $\tilde{\cC}_E\setminus(X\cup Y)$ 之中，因此可以通过
\[
\begin{array}{rcl}
s_j:\ E\times D & \to & \cC\\
(t,\tau) & \mapsto & (s_j(t),\tau)\in Z
\end{array}
\]
定义全纯截面 $s_j$。

用同样的方式可以定义全纯坐标 $\eta_j$。容易证明 $(\pi:\cC\to E\times D;\,s_1,\dots,s_N)$（相应地 $(\pi:\cC\to E\times D;\,s_1,\dots,s_N;\eta_1,\dots,\eta_N)$）同构于我们原来的族 $\mathfrak{F}^{(0)}$（相应地 $\mathfrak{F}$）。因此，以下我们把 $\mathfrak{F}^{(0)}$ 与 $\mathfrak{F}$ 和上述构造的族等同起来。

对每个点 $(u,\tau)\in E^{(1)}\times D$，置
\[
\begin{array}{rcl}
\cC_{(u,\tau)} & = & \pi^{-1}((u,\tau))\\[2pt]
U_\alpha(u,\tau) & = & U_\alpha\cap\cC_{u,\tau},\qquad 3\le\alpha\le A\\[2pt]
U_1(u,\tau) & = & S\cap Z\cap\cC_{u,\tau}\\[2pt]
U_2(u,\tau) & = & S\cap W\cap\cC_{u,\tau}.
\end{array}
\]
于是，对每个 $\tau\ne0$，$\mathcal U(u,\tau)=\{U_\alpha(u,\tau)\}_{1\le\alpha\le A}$ 是曲线 $\cC_{(u,\tau)}$ 的开覆盖。

\begin{lemma}[5.16]
对每个点 $(u,\tau)\in E\times D$（$\tau\ne0$），向量场 $\tau\dfrac{\partial}{\partial\tau}$ 在小平–Spencer 映射
\[
\rho: T_{(u,\tau)}(E\times D)\to H^1(\cC_{(u,\tau)},\Theta_{\cC_{(u,\tau)}})
\]
下的像 $\rho\bigl(\tau\frac{\partial}{\partial\tau}\bigr)$，由关于上述覆盖 $\mathcal U(u,\tau)$ 的 Čech 上同调类 $\{\theta_{\alpha\beta}(u,\tau)\}\in\check H^1(\mathcal U(u,\tau),\Theta_{\cC_{(u,\tau)}})$ 给出，其中
\[
\begin{array}{rcl}
\theta_{12}(u,\tau) & = & z\dfrac{\partial}{\partial z}\\[6pt]
\theta_{21}(u,\tau) & = & -\theta_{12}(u,\tau)\\[6pt]
\theta_{\alpha\beta}(u,\tau) & = & 0\qquad\text{若 } (\alpha,\beta)\ne(1,2)\text{ 或 }(2,1).
\end{array}
\]
\end{lemma}

\noindent\textit{证明.}\quad 由上述等价关系，在 $U_1(u,\tau)\cap U_2(u,\tau)$ 上有
\[
z=\frac{\tau}{w}.
\]
若 $U_\alpha(u,\tau)\cap U_\beta(u,\tau)\ne\emptyset$ 且 $(\alpha,\beta)$ 既非 $(1,2)$ 也非 $(2,1)$，则 $U_\alpha(u,\tau)$ 与 $U_\beta(u,\tau)$ 的局部坐标之间的关系不依赖于 $\tau$。因此，由小平–Spencer 映射的定义（见注 1.37）得到
\[
\begin{array}{rcl}
\rho\bigl(\tau\dfrac{\partial}{\partial\tau}\bigr)_{12} & = & \dfrac{\tau}{w}\dfrac{\partial}{\partial\tau}
   = z\dfrac{\partial}{\partial z}\\[10pt]
\rho\bigl(\tau\dfrac{\partial}{\partial\tau}\bigr)_{21} & = & w\dfrac{\partial}{\partial w}
   = -z\dfrac{\partial}{\partial z}\\[10pt]
\rho\bigl(\tau\dfrac{\partial}{\partial\tau}\bigr)_{\alpha\beta} & = & 0\qquad
   \text{若 } (\alpha,\beta)\ne(1,2)\text{ 或 }(2,1).
\end{array}
\]
\hfill$\square$

% [Tutor Note: 原文 $\rho(\tau\partial/\partial\tau)_{12}=\frac{\tau}{w}\frac{\partial}{\partial\tau}$
%  中的 $\partial/\partial\tau$ 应为 $\partial/\partial z$：Kodaira–Spencer 类取的是转移函数
%  $z=\tau/w$ 对 $\tau$ 求导后乘 $\partial/\partial z$，即 $\frac{\partial z}{\partial\tau}\frac{\partial}{\partial z}
%  =\frac{1}{w}\frac{\partial}{\partial z}$，再乘 $\tau$ 得 $\frac{\tau}{w}\frac{\partial}{\partial z}=z\frac{\partial}{\partial z}$，
%  与末式一致；若真写成 $\partial/\partial\tau$，它就不是曲线上的向量场。]

% [原文 p.106]
它在 (5.23) 中定义。由于 $\ell_j(\xi_j)\dfrac{d}{d\xi_j}\in\cO_E((\xi_j))$，我们可以把 $\vec\ell$ 视为 $\mathfrak{F}^{(0)}$ 上的 $N$-元形式向量场，即 $\ell_j(\xi_j)\dfrac{d}{d\xi_j}\in\cO_{E\times D}((\xi_j))$。

\begin{corollary}[5.17]
在 $B=E\times D$ 上有
\[
\theta(\vec\ell)=\tau\frac{\partial}{\partial\tau},
\]
其中映射 $\theta$ 在 (4.15) 中定义。
\end{corollary}

\noindent\textit{证明.}\quad 由于推论中等式的两边都是 $B$ 上的全纯向量场，只需对 $\tau\ne0$ 证明该等式。

考虑正合列
\[
0\to\Theta_{\cC/B}(-S)\to\Theta_{\cC/B}((m-1)S)\to
\bigoplus_{j=1}^{N}\bigoplus_{k=1}^{m}\cO_B\xi_j^{-m+k}\frac{d}{d\xi_j}\to 0
\]
（$m$ 是充分大的正整数）。$\vec\ell$ 定义了该正合列第三项中的一个元素 $\vec\ell$。另一方面，对每个 $(u,\tau)\in E\times D$（$\tau\ne0$），$\tilde{\cC}_E$ 上的亚纯向量场 $\tilde\ell$ 在 $\cC_{u,\tau}\setminus\{U_2(u,\tau)\setminus(U_1(u,\tau)\cap U_2(u,\tau))\}$ 上定义了亚纯向量场 $\tilde\ell_{u,\tau}$，并在 $U_2(u,\tau)$ 上定义了 $\tilde\ell'_{u,\tau}=\dfrac12 w\dfrac{\partial}{\partial w}$，使得这两个向量场在上述正合列中具有相同的像 $\vec\ell$。因此 $\vec\ell$ 在映射
\[
\bigoplus_{j=1}^{N}\bigoplus_{k=1}^{m}\cO_B\xi_j^{-m+k}\frac{d}{d\xi_j}
\longrightarrow R^1\pi_*(\Theta_{\cC/B}(-S))
\]
下的像，在点 $(u,\tau)$ 处由元素
\[
\{\theta_{\alpha,\beta}(u,\tau)\}\in H^1(\cC_{(u,\tau)},\Theta_{\cC_{(u,\tau)}})
\]
给出，其中在 $U_1(u,\tau)\cap U_2(u,\tau)$ 上有
\[
\begin{array}{rcl}
\theta_{12}(u,\tau) & = & \tilde\ell_{u,\tau}\big|_{U_1(u,\tau)}
   -\tilde\ell'_{u,\tau}\big|_{U_2(u,\tau)}\\[6pt]
 & = & \dfrac12 z\dfrac{\partial}{\partial z}-\dfrac12 w\dfrac{\partial}{\partial w}\\[6pt]
 & = & z\dfrac{\partial}{\partial z},
\end{array}
\]
\[
\theta_{21}(u,\tau)=-\theta_{12}(u,\tau),
\]
而在 $U_\alpha(u,\tau)\cap U_\beta(u,\tau)$ 上（$(\alpha,\beta)\ne(1,2),(2,1)$）有
\[
\theta_{\alpha\beta}(u,\tau)=0.
\]
于是 $\vec\ell$ 定义了引理 5.16 中给出的上同调类。因此对 $\tau\ne0$ 等式成立。\hfill$\square$

设 $\ell$ 是亚纯向量场，其 $\theta(\ell)$ 由引理 5.15 给出，而 $\ell_j$ 是 $\tilde\ell$ 在 $Q_j$ 处的形式 Laurent 展开。如 (5.23) 那样置 $\vec\ell=(\ell_1,\dots,\ell_N)$。设 $\langle\tilde\Psi|\in\cH^\dagger_\lambda(E)[[\tau]]$ 是 (4.26) 中定义的元素。置
\[
\langle\hat\Psi|=\tau^{\Delta_\mu}\langle\tilde\Psi|.
\]

% [Tutor Note: 原文"其 $\theta(\ell)$ 由引理 5.15 给出"疑为"由推论 5.17 给出"：
%  $\theta(\vec\ell)=\tau\partial/\partial\tau$ 是推论 5.17 的结论，引理 5.15 里没有 $\theta$。
%  同句前面用 $\ell$、后面用 $\tilde\ell$ 混用同一对象，亦属原文笔误（照排）。]

% [原文 p.107]
\figomit{5.2}{曲线 $C'_{(u,\tau)}$ 的边界。}

我们将证明 $\langle\hat\Psi|$ 收敛，且是边界附近 $\cV^\dagger_\lambda(\mathfrak{F})$ 的多值全纯截面；办法是证明 $\langle\hat\Psi|$ 是 Fuchs 型微分方程的形式解。

\begin{theorem}[5.18]
形式幂级数 $\langle\hat\Psi|$ 是 Fuchs 型微分方程
\begin{equation}
\Bigl(\tau\frac{d}{d\tau}-T[\vec\ell]+a(\vec\ell)\Bigr)\langle\hat\Psi|=0 \tag{5.24}
\end{equation}
的形式解。此外，$\langle\hat\Psi|$ 收敛，并在边界附近定义 $\cV^\dagger_\lambda(\mathfrak{F})$ 的一个多值全纯截面。
\end{theorem}

\noindent\textit{证明.}\quad 我们可以假设 $B^{(0)}$ 足够小，使得存在对称双微分
\[
\omega\in H^0\bigl(\cC^{(0)}\otimes_{B^{(0)}}\cC^{(0)},\ \omega^{\otimes2}_{\cC^{(0)}/B^{(0)}}(2\Delta)\bigr),
\]
且 $\operatorname{Res}^2(\omega)=1$。由 (5.2) 我们有
\[
\langle\hat\Psi|\tilde T(u)|\Phi\rangle du^2\in H^0\bigl(\cC,\ \omega^{\otimes2}_{\cC/B}(*S^{(1)})\bigr).
\]

设 $\tilde\ell=\ell(z)\dfrac{\partial}{\partial z}$ 是引理 5.15 中给出的亚纯向量场。于是对 $(u,\tau)\in E\times D$（$\tau\ne0$），
\[
\ell(z)\langle\hat\Psi|\tilde T(z)|\Phi\rangle dz
\]
是
\[
C'_{(u,\tau)}=\cC_{(u,\tau)}\setminus\bigl\{(x,y,\tau)\in S_0\ \big|\ |x|\le\eps\ \text{或}\ |y|\le\eps\bigr\}
\]
上的亚纯形式，其中 $\eps<1$ 是充分小的正数。

$C'_{(u,\tau)}$ 的边界由两条互不相交的简单闭曲线 $\gamma_+,\gamma_-$ 组成。我们如此选取 $\gamma_\pm$ 的定向，使得 $C'_{u,\tau}$ 位于 $\gamma_\pm$ 的右手侧。于是由定理 3.26、(5.2) 与 (5.3)，我们有
\[
\begin{array}{rcl}
&&\displaystyle\frac{1}{2\pi\sqrt{-1}}\int_{\gamma_+}\ell(z)\langle\hat\Psi|\tilde T(z)|\Phi\rangle dz
+\frac{1}{2\pi\sqrt{-1}}\int_{\gamma_-}\ell(w)\langle\hat\Psi|\tilde T(w)|\Phi\rangle dw\\[10pt]
&=& \displaystyle\sum_{j=1}^{N}\operatorname{Res}_{Q_j}\bigl(\ell(u)\langle\hat\Psi|\tilde T(u)|\Phi\rangle du\bigr)\\[10pt]
&=& \displaystyle\sum_{j=1}^{N}\bigl\langle\hat\Psi\big|\rho_j\bigl(\operatorname{Res}_{\xi_j=0}
   (\ell_j(\xi_j)T(\xi_j)d\xi_j)\bigr)\big|\Phi\bigr\rangle
-\frac{c_v}{12}\sum_{j=1}^{N}\operatorname{Res}_{\xi_j=0}\bigl(\ell_j(\xi_j)S_\omega(\xi_j)\bigr)
   \langle\hat\Psi|\Phi\rangle\\[10pt]
&=& \displaystyle\sum_{j=1}^{N}\bigl\langle\hat\Psi\big|\rho_j(T[\ell_j])\big|\Phi\bigr\rangle
-a(\vec\ell)\langle\hat\Psi|\Phi\rangle.
\end{array}
\]
% [原文 p.108]
另一方面，在 $\gamma_+$ 上有 $\ell(z)\dfrac{d}{dz}=\dfrac12 z\dfrac{d}{dz}$。因此由 (5.2) 有
\[
\begin{array}{rcl}
&&\displaystyle\frac{1}{2\pi\sqrt{-1}}\int_{\gamma_+}\ell(z)\langle\hat\Psi|\tilde T(z)|\Phi\rangle dz\\[12pt]
&=& \displaystyle\frac{1}{4\pi\sqrt{-1}}\int_{\gamma_+}
\Bigl(\ell(z)\langle\hat\Psi|T(z)|\Phi\rangle-\frac{c_v}{12}\ell(z)S_\omega(z)\langle\hat\Psi|\Phi\rangle\Bigr)dz\\[12pt]
&=& \displaystyle\frac{1}{4\pi\sqrt{-1}}\int_{\gamma_+}\ell(z)\langle\hat\Psi|T(z)|\Phi\rangle dz,
\end{array}
\]
因为 $S_\omega(z)dz^2$ 在 $z=0$ 处全纯。因此由 (4.24) 与 (4.25) 以及定理 3.26 的 (1)，我们有
\[
\begin{array}{rcl}
&&\displaystyle\frac{1}{2\pi\sqrt{-1}}\int_{\gamma_+}\ell(z)\langle\hat\Psi|\tilde T(z)|\Phi\rangle dz\\[12pt]
&=& \displaystyle\frac12\sum_{d=0}^{\infty}\tau^{\Delta_\mu+d}\sum_{i=1}^{m_d}
\int_{\gamma_+}z\,\langle\Psi_d|T(z)|v_i(d)\otimes v^i(d)\otimes\Phi\rangle dz\\[12pt]
&=& \displaystyle\frac12\sum_{d=0}^{\infty}\tau^{\Delta_\mu+d}\sum_{i=1}^{m_d}
\langle\Psi_d|L_0(v_i(d))\otimes v^i(d)\otimes\Phi\rangle\\[12pt]
&=& \displaystyle\frac12\sum_{d=0}^{\infty}\sum_{i=1}^{m_d}
(\Delta_\mu+d)\tau^{\Delta_\mu+d}\langle\Psi_d|v_i(d)\otimes v^i(d)\otimes\Phi\rangle.
\end{array}
\]
同样地有
\[
\begin{array}{rcl}
&&\displaystyle\frac{1}{2\pi\sqrt{-1}}\int_{\gamma_-}\ell(w)\langle\hat\Psi|\tilde T(w)|\Phi\rangle dw\\[12pt]
&=& \displaystyle\frac12\sum_{d=0}^{\infty}\sum_{i=1}^{m_d}
(\Delta_{\mu^\dagger}+d)\tau^{\Delta_{\mu^\dagger}+d}
\langle\Psi_d|v_i(d)\otimes v^i(d)\otimes\Phi\rangle.
\end{array}
\]
由于 $\Delta_\mu=\Delta_{\mu^\dagger}$，我们得到
\[
\begin{array}{rcl}
&&\displaystyle\sum_{j=1}^{N}\langle\hat\Psi|\rho_j(T[\ell_j])|\Phi\rangle
-a(\vec\ell)\langle\hat\Psi|\Phi\rangle\\[12pt]
&=& \displaystyle\sum_{d=0}^{\infty}\sum_{i=1}^{m_d}
(\Delta_\mu+d)\tau^{\Delta_\mu+d}\langle\Psi_d|v_i(d)\otimes v^i(d)\otimes\Phi\rangle.
\end{array}
\]
另一方面，我们有
\[
\{\langle\hat\Psi|\tau\frac{d}{d\tau}\}|\Phi\rangle
=\sum_{d=0}^{\infty}\sum_{i=1}^{m_d}(\Delta_\mu+d)\tau^{\Delta_\mu+d}
\langle\Psi_d|v_i(d)\otimes v^i(d)\otimes\Phi\rangle.
\]
因此 $\langle\hat\Psi|$ 是微分方程
\[
\Bigl(\tau\frac{d}{d\tau}\langle\hat\Psi|\Bigr)-\bigl(\langle\hat\Psi|T[\vec\ell]\bigr)
+a(\vec\ell)\langle\hat\Psi|=0
\]
的形式解。

最后我们证明 $\langle\hat\Psi|$ 收敛，并在边界附近定义 $\cV^\dagger_\lambda(\mathfrak{F})$ 的多值截面。由定理 4.19，$\cV^\dagger_\lambda(\mathfrak{F})$ 是局部自由的。取局部平凡化
\[
p:\cV^\dagger_\lambda(\mathfrak{F})\simeq\cO_B^{\oplus n}.
\]

% [原文 p.109]
于是，对从左侧作用在 $\cV^\dagger_\lambda(\mathfrak{F})$ 上的微分算子 $D:=\tau\dfrac{d}{d\tau}-T[\vec\ell]+a([\vec\ell])$，作用在 $\cO_B^{\oplus n}$ 上的算子 $P:=p\circ D\circ p^{-1}$ 是一阶 Fuchs 型微分方程组。另一方面，由引理 4.18，局部平凡化 $p$ 诱导出层 $\cV^\dagger_{\vec\lambda/E}$ 的局部平凡化
\[
\hat p:\hat{\cV}^\dagger_{\vec\lambda/E}\simeq\hat\cO^{\oplus n}_{\cB/E}.
\]
于是形式解 $\langle\hat\Psi|$ 可以看作 $n$ 元形式幂级数，它是 Fuchs 型微分方程组 $P$ 的解。由 Fuchs 型微分方程的理论，形式解总是收敛的。因此 $\langle\hat\Psi|$ 在边界附近确定 $\cV^\dagger_\lambda(\mathfrak{F})$ 的一个多值截面。\hfill$\square$

\subsection{单点椭圆曲线的共形块}

让我们用上一节的论证来计算带点椭圆曲线的共形块。

上半平面 $H$ 上椭圆曲线的通用族 $\pi:\mathcal E\to H$ 已在例 1.1 中构造。映射 $s:H\ni\tau\mapsto(\tau,[0])$ 与 $\C$ 的坐标 $\xi$ 给出了 $H$ 上带坐标的单点椭圆曲线族 $\mathfrak{F}=(\pi:\mathcal E\to H;s;\xi)$。点 $\tau\in H$ 上的纤维 $\pi^{-1}(\tau)$ 是椭圆曲线 $E_\tau=\C/(1,\tau)$，其周期矩阵为 $(1,\tau)$。为构造亏格一的稳定单点曲线族及其坐标，我们需要在去心单位圆盘 $D^*=\{z\in\C\mid |z|<1\}$ 上构造椭圆曲线族。置
\[
q=e^{2\pi\sqrt{-1}\tau}.
\]
于是椭圆曲线 $E_\tau$ 可以表示成商流形 $\C^*/\langle q\rangle$，其中 $\C^*=\C\setminus\{0\}$ 由自同构 $w\mapsto qw$ 生成的无限循环群作用，这里 $w=e^{2\pi\sqrt{-1}z}$。于是我们有通用椭圆曲线族 $\hat\pi:\hat{\mathcal E}^*\to D^*$，使得 $q\in D^*$ 上的纤维 $\hat\pi^{-1}(q)$ 是椭圆曲线 $\C^*/\langle q\rangle$。众所周知，该族可延拓为单位圆盘 $D$ 上的族 $\hat\pi:\hat{\mathcal E}\to D$，使得原点上的纤维是带一个通常二重点的有理曲线。该族具有全纯截面 $\hat s:D\to\hat{\mathcal E}$，它在 $D^*$ 上的限制由 $\hat s(q)=[1]$ 给出，其中 $[1]$ 是椭圆曲线 $\C^*/\langle q\rangle$ 的原点。于是 $\hat{\mathfrak{F}}=(\hat\pi:\hat{\mathcal E}\to D;\hat s;w-1)$ 是亏格一的稳定单点曲线族及其坐标。

% [Tutor Note: 原文说 $D^*$ 是"去心单位圆盘"，却写成 $D^*=\{z\in\C\mid |z|<1\}$，漏了 $0<|z|$
%  （或应写成 $\{q\in\C\mid 0<|q|<1\}$）；否则 $q=0$ 落在 $D^*$ 里，而 $\C^*/\langle q\rangle$
%  在 $q=0$ 处退化，与"去心"矛盾。]

设
\[
\nu:\tilde{\mathcal E}(0)=\PP^1\to\hat{\mathcal E}(0)
\]
是曲线 $\hat{\mathcal E}(0)=\hat\pi^{-1}(0)$ 的正规化。我们这样选取 $\PP^1$ 的坐标，使得
\[
\begin{array}{rcl}
\nu(0) & = & \nu(\infty)=\text{二重点},\\[2pt]
\nu(1) & = & s(0).
\end{array}
\]

% [Tutor Note: 原文写 $\nu:\tilde{\mathcal E}(0)=\PP^1\to\hat{\mathcal E}(0)$，其中 $\tilde{\mathcal E}(0)$
%  应为 $\hat{\mathcal E}(0)$：正规化的定义域是纤维 $\hat{\mathcal E}(0)=\hat\pi^{-1}(0)$ 的正规化
%  $\PP^1$，而 $\tilde{\mathcal E}$ 此前并未引入。照排并在此指出。]

我们知道 $\mathcal V^\dagger_0((\mathcal E(0);\hat s(0),u))$ 典范同构于
\[
\bigoplus_{\lambda\in P_\ell}\cV^\dagger_{\lambda^\dagger,\lambda,0}
\bigl((\PP^1;\infty,0,1;\eta,\xi,\xi-1)\bigr).
\]
为简单起见，置
\[
V^\dagger_{\lambda^\dagger,\lambda,0}
=\cV^\dagger_{\lambda^\dagger,\lambda,0}\bigl((\PP^1;\infty,1,0;\eta,\xi-1,\xi)\bigr).
\]

% [原文 p.110]
向量空间 $V^\dagger_{\lambda^\dagger,\lambda,0}$ 是一维的，且存在元素 $\langle\Psi|\in V^\dagger_{\lambda^\dagger,\lambda,0}$，使得
\[
\langle\Psi|v^\dagger\otimes v\otimes 0\rangle=(v^\dagger|v)
\]
对任意 $|v^\dagger\rangle\in\cH_{\lambda^\dagger}$ 与 $|v\rangle\in\cH_\lambda$ 成立，其中 $|0\rangle$ 是 $\cH_0$ 的最高权向量，而 $(\ |\ )$ 是引理 2.17 中定义的双线性配对
\[
\cH_{\lambda^\dagger}\otimes\cH_\lambda\to\C.
\]
让我们像上一节那样，从 $\langle\Psi|$ 构造元素
\[
\langle\tilde\Psi^{(\lambda)}|\in\cV^\dagger_0(\mathfrak{F}).
\]
对任意非负整数 $d$，设 $\{v^k(d)\}$ 与 $\{v_k(d)\}$（$k=1,2,\dots,m_d$）分别是 $\cH_{\lambda^\dagger,d}$ 与 $\cH_{\lambda,d}$ 的基，使得
\[
(v^j(d)|v_k(d))=\delta^j_k.
\]
对每个元素 $|\phi\rangle\in\cH_0$，由下式定义元素 $\langle\tilde\Psi_d|\in\cH_0^\dagger$：
\[
\langle\tilde\Psi_d|\phi\rangle=\sum_{i=1}^{m_d}
\langle\tilde\Psi|v^i(d)\otimes0\otimes v_i(d)\otimes\phi\rangle.
\]
现在由下式定义 $\langle\tilde\Psi^{(\lambda)}|$：
\[
\langle\tilde\Psi^{(\lambda)}|\phi\rangle=\sum_{d=0}^{\infty}\langle\tilde\Psi_d|\phi\rangle q^d.
\]
这是关于 $q$ 的形式幂级数，且由定理 5.18 它收敛。若置 $|\phi\rangle=|0\rangle$，则
\begin{equation}
\langle\hat\Psi^{(\lambda)}|0\rangle=q^{\Delta_\lambda}\langle\tilde\Psi^{(\lambda)}|0\rangle
=\sum_{d=0}^{\infty}\sum_{i=1}^{m_d}(v^i(d)|v_i(d))q^d
=\sum_{d=0}^{\infty}\dim\cH_{\lambda,d}q^{\Delta_\lambda+d}. \tag{5.25}
\end{equation}
由于 $L_0$ 在 $\cH_\lambda(d)$ 上以 $\Delta_\lambda+d$ 作乘法，幂级数 (5.25) 常写成
\[
\Tr_{\cH_\lambda}q^{L_0}
\]
的形式。

考虑乘以 $q^{-c_v/24}$ 之后的幂级数
\begin{equation}
\Tr_{\cH_\lambda}q^{L_0-c_v/24}
=\sum_{d=0}^{\infty}\dim\cH_{\lambda,d}q^{\Delta_\lambda+d-c_v/24}, \tag{5.26}
\end{equation}
它比原来的幂级数 (5.25) 更好，因为这对应于最高权模 $\cH_\lambda$ 的特征标 $\chi_\lambda$。可积 $\hat{\mathfrak g}$-模 $\cH_\lambda$ 的特征标可以用 theta 函数描述如下：
\begin{equation}
\chi_\lambda(z,\tau)=
\frac{\sum_{w\in W}(-1)^{|W|}\Theta_{w(\lambda+\rho),g^*+\ell}}
     {\sum_{w\in W}(-1)^{|W|}\Theta_{w(\rho),g^*}} \tag{5.27}
\end{equation}
其中 $W$ 是李代数 $\mathfrak g$ 的 Weyl 群，且
\begin{equation}
\Theta_{\lambda,\ell}(z,\tau)=
\sum_{\mu\in\ell Q^\vee+\lambda}
e^{\pi\sqrt{-1}\tau(\mu,\mu)/\ell+2\pi\sqrt{-1}(z,\mu)}, \tag{5.28}
\end{equation}
其中 $Q^\vee=\sum\Z\alpha_i^\vee$ 是李代数 $\mathfrak g$ 的余根格。

它是 $D^*$ 上的多值全纯函数，但在上半平面 $H$ 上是单值的。设 $\tilde{\mathfrak F}$ 是族 $\mathfrak{F}|_{D^*}$ 到 $H$ 的拉回。则 $\chi_\lambda(z,\tau)$ 是

% [Tutor Note: (5.25) 的第二个等号漏了因子 $q^{\Delta_\lambda}$：由 $\langle\tilde\Psi^{(\lambda)}|$ 的
%  定义，$\langle\tilde\Psi^{(\lambda)}|0\rangle=\sum_d\sum_i(v^i(d)|v_i(d))q^d$，
%  所以 $q^{\Delta_\lambda}\langle\tilde\Psi^{(\lambda)}|0\rangle
%  =q^{\Delta_\lambda}\sum_d\sum_i(v^i(d)|v_i(d))q^d=\sum_d\dim\cH_{\lambda,d}q^{\Delta_\lambda+d}$。
%  按原书照排，正确中间项应带上 $q^{\Delta_\lambda}$。]

% [Tutor Note: (5.27) 的符号 $(-1)^{|W|}$ 照原文排；通常写作 $(-1)^{|w|}$ 或 $\det(w)$、
%  $(-1)^{\ell(w)}$（$\ell(w)$ 为 $w$ 的长度）。$|W|$ 是 Weyl 群的阶，与 $w$ 无关，
%  若真按 $(-1)^{|W|}$ 理解则求和号内出现整体常数符号，故此处显系 $|w|$ 之误。]

% [原文 p.111]
是 $H$ 上 $\cV_0^\dagger(\tilde{\mathfrak F})$ 的全纯截面。于是由定理 5.18 得到如下定理。

\begin{theorem}[5.19]
带标号 $0$ 的单点椭圆曲线的共形块的一组基由所有特征标 $\chi_\lambda$（$\lambda\in P_\ell$）给出。此外，$\{\chi_\lambda\}_{\lambda\in P_\ell}$ 是共形块丛 $\cV_0^\dagger(\tilde{\mathfrak F})$ 的平坦截面的一组基。
\end{theorem}

现在，如例 1.2 所示，群 $\SL(2,\Z)$ 全纯地作用在上半平面 $H$ 上的族 $\mathfrak{F}$ 上，从而也等变地作用在共形块层 $\cV_0^\dagger(\mathfrak{F})\to H$ 上。群 $\SL(2,\Z)$ 由 $H$ 上的两个模变换生成：
\[
T=\begin{pmatrix}1&1\\0&1\end{pmatrix},\qquad
S=\begin{pmatrix}0&-1\\1&0\end{pmatrix},
\]
它们作用在族上为
\[
T:(\tau,z)\mapsto(\tau+1,z),\qquad
S:(\tau,z)\mapsto\Bigl(-\frac{1}{\tau},\frac{z}{\tau}\Bigr).
\]
于是显然有
\[
\chi_\lambda(z,\tau+1)=q^{\Delta_\lambda-c_v/24}\chi_\lambda(z,\tau).
\]
又已知
\[
\chi_\lambda\Bigl(\frac{z}{\tau},-\frac{1}{\tau}\Bigr)
=\sum_{\mu\in P}S_{\lambda,\mu}\chi_\mu(z,\tau),
\]
其中 $S_{\lambda,\mu}$ 由下式给出
\begin{equation}
S_{\lambda,\mu}=S_{\lambda,0}\,\mathrm{ch}_\mu
\Bigl(e^{-\frac{2\pi\sqrt{-1}(\lambda+\rho)}{\ell+g^*}}\Bigr), \tag{5.29}
\end{equation}
\begin{equation}
S_{\lambda,0}=|P/(\ell+g^*)Q^\vee|^{-1/2}
\prod_{\alpha>0}2\sin\Bigl(\frac{\pi(\lambda+\rho,\alpha)}{\ell+g^*}\Bigr), \tag{5.30}
\end{equation}
这里 $\mathrm{ch}_\mu$ 是单李代数 $\mathfrak g$ 的权为 $\mu$ 的特征标。（例如见 [Kac]，§13 定理 13.9。）

矩阵 $(S_{\lambda,\mu})$ 称为 $S$-矩阵，它在 Verlinde 公式中起重要作用。

% [Tutor Note: $S$-变换的求和指标原文写作 $\mu\in P$，但 $\chi_\mu$ 只对 $\mu\in P_\ell$
%  定义（定理 5.19），故应为 $\mu\in P_\ell$；同理 (5.30) 中的 $P$ 指权格。]

\subsection{Verlinde 公式}

让我们固定复单李代数 $\mathfrak g$ 与如上所述的水平 $\ell$。设 $\N^{(P_\ell)}$ 是由 $P_\ell$ 生成的交换幺半群，即
\[
\N^{(P_\ell)}=\Bigl\{\sum_{\lambda\in P_\ell}n_\lambda\lambda
\ \Big|\ n_\lambda\ \text{为非负整数}\Bigr\}.
\]
对任意非负整数 $g$，我们定义映射
\[
N_g:\N^{(P_\ell)}\to\Z
\]
如下：
\[
N_g(\lambda_1+\dots+\lambda_m)=\dim_{\C}\cV^\dagger_{\vec\lambda}
\bigl((C;Q_1,\dots,Q_m;\xi_1,\dots,\xi_m)\bigr),
\]
其中 $(C;Q_1,\dots,Q_m;\xi_1,\dots,\xi_m)$ 是亏格 $g$ 的 $m$ 点黎曼曲面，而 $\vec\lambda=(\lambda_1,\dots,\lambda_m)$。
对 $\alpha,\beta,\gamma\in P_\ell$，如下定义 $N^\gamma_{\alpha,\beta}$：
\[
N^\gamma_{\alpha,\beta}=\dim_{\C}\cV^\dagger_{\alpha,\beta,\gamma^\dagger}
(\PP^1;Q_1,Q_2,Q_3;\xi_1,\xi_2,\xi_3).
\]

% [原文 p.112]
注意我们在点 $Q_3$ 处使用标号 $\gamma^\dagger$。按定义有
\begin{equation}
N^\gamma_{\alpha,\beta}=N^\gamma_{\beta,\alpha}. \tag{5.31}
\end{equation}
置
\[
\Z^{(P_\ell)}=\bigoplus_{\lambda\in P_\ell}\Z\lambda.
\]
在 $\Z^{(P_\ell)}$ 上如下定义乘法 "$\cdot$"：
\begin{equation}
\alpha\cdot\beta=\sum_{\lambda\in P_\ell}N^\lambda_{\alpha,\beta}\,\lambda, \tag{5.32}
\end{equation}
其中 $\alpha,\beta\in P_\ell$，并把它 $\Z$-双线性地延拓。

\begin{proposition}[5.20]
乘法 (5.32) 在 $\Z^{(P_\ell)}$ 上定义了交换环结构。该环称为 Verlinde 代数。
\end{proposition}

\noindent\textit{证明.}\quad 由 (5.31)，乘法是交换的。按定义有
\[
\begin{array}{rcl}
(\alpha\cdot\beta)\cdot\gamma & = &
\displaystyle\sum_{\lambda\in P_\ell}N^\lambda_{\alpha\beta}\,\lambda\cdot\gamma
=\sum_{\lambda,\mu\in P_\ell}N^\lambda_{\alpha\beta}N^\mu_{\lambda\gamma}\,\mu,\\[12pt]
\alpha\cdot(\beta\cdot\gamma) & = &
\displaystyle\sum_{\lambda\in P_\ell}N^\lambda_{\beta\gamma}\,\alpha\cdot\lambda
=\sum_{\lambda,\mu\in P_\ell}N^\lambda_{\beta\gamma}N^\mu_{\alpha\lambda}\,\mu.
\end{array}
\]
由因子化（定理 3.19）我们有
\[
\sum_{\lambda\in P_\ell}N^\lambda_{\alpha\beta}N^\mu_{\lambda\gamma}
=\sum_{\lambda\in P_\ell}N^\lambda_{\beta\gamma}N^\mu_{\alpha\lambda}.
\]
因此乘法是结合的。\hfill$\square$

注意 $0\in P_\ell$ 并不是 Verlinde 代数的零元。还要注意，对合 $\dagger$ 可以 $\Z$-线性地延拓为 $\Z^{(P_\ell)}$ 上的一个对合。

让我们固定同构
\[
\lambda:\{0,1,\dots,m\}\to P_\ell,
\]
使得
\[
\lambda(0)=0\in P_\ell.
\]
在集合 $\{0,1,\dots,m\}$ 上引入对合 $\dagger$，使得
\[
\lambda(i^\dagger)=\lambda(i)^\dagger.
\]
置
\begin{equation}
N^k_{i,j}=N^{\lambda(k)}_{\lambda(i),\lambda(j)}. \tag{5.33}
\end{equation}
设 $N_i$ 是以 $N^k_{i,j}$ 为 $(j,k)$-元素的矩阵。对 $n\times n$ 矩阵 $M$，我们用 $M^k_j$ 表示 $M$ 的 $(j,k)$-元素。于是由因子化定理（定理 3.19）可以证明如下定理。

\begin{theorem}[5.21]
对 $N$ 点 $\PP^1$，$\mathfrak X=(\PP^1;a_1,\dots,a_N;z-a_1,\dots,z-a_N)$，我们有
\[
\dim_{\C}\cV^\dagger_{\lambda(i_1),\dots,\lambda(i_N)}(\mathfrak X)
=(N_{i_1}N_{i_2}\cdots N_{i_N})^0_0.
\]
对亏格 $g$ 的 $N$ 点曲线 $\mathfrak X=(C;Q_1,\dots,Q_N;\xi_1,\dots,\xi_N)$（带坐标），我们有
\[
\dim_{\C}\cV^\dagger_{\lambda(i_1),\dots,\lambda(i_N)}(\mathfrak X)
=(N_{i_1}N_{i_2}\cdots N_{i_N}W^g)^0_0.
\]
\end{theorem}

% [Tutor Note: 定理 5.21 第二式中的 $W$ 在此处尚未定义（它在后文才作为作用在
%  Verlinde 代数上的算子出现），这是本节的一处跳步。可先核对第一式：由 (5.33)，
%  $N_i$ 的第 $(j,k)$ 元是 $N^{\lambda(k)}_{\lambda(i),\lambda(j)}$，即三点函数；
%  连乘 $N_{i_1}\cdots N_{i_N}$ 再取 $(0,0)$ 元，正是因子化定理把 $N$ 点共形块
%  逐次粘合后剩下的 $0$ 标号（真空）系数，故等于共形块的维数。]

% [原文 p.113]
其中
\begin{equation}
W=\sum_{k=0}^{m}(\Tr N_{k^\dagger})N_k. \tag{5.34}
\end{equation}
此外，若 $g\ge1$，则维数等于
\[
\Tr(N_{i_1}N_{i_2}\cdots N_{i_N}W^{g-1}).
\]

\noindent\textit{证明.}\quad 对亏格 $0$ 的情形，在 $N$ 点 $\PP^1$ 上添加两点 $Q_0$ 与 $Q_{N+1}$，得到
\[
\mathfrak X'=(\PP^1;a_0,a_1,\dots,a_N,a_{N+1};z-a_0,\dots,z-a_{N+1}).
\]
于是由真空的传播（定理 3.15）有
\[
\dim_{\C}\cV^\dagger_{\lambda(i_1),\dots,\lambda(i_N)}(\mathfrak X)
=\dim_{\C}\cV^\dagger_{0,\lambda(i_1),\dots,\lambda(i_N),0}(\mathfrak X').
\]
现在我们可以利用把 $(N+2)$ 点 $\PP^1$ 退化成一个由 $N$ 个 $\PP^1$ 组成的带点亏格 $0$ 稳定曲线来证明定理。

\figomitn{把 $(N+2)$ 点 $\PP^1$ 退化所得、由 $N$ 个 $\PP^1$ 组成的亏格 $0$ 带点稳定曲线；标记依次为 $\lambda(i_1),\dots,\lambda(i_N)$，两端为真空标号 $0$。}

由因子化定理与 (5.31) 我们有
\[
\begin{array}{rcl}
\dim_{\C}\cV^\dagger_{0,\lambda(i_1),\dots,\lambda(i_n),0}(\mathfrak X')
&=& \displaystyle\sum_{k_1,\dots,k_N}N^{k_1}_{i_1,0}N^{k_2}_{k_1,i_2}\cdots N^{0}_{k_N,i_n}\\[10pt]
&=& \displaystyle\sum_{k_1,\dots,k_N}N^{k_1}_{i_1,0}N^{k_2}_{i_2,k_1}\cdots N^{0}_{i_n,k_N}
=(N_{i_1}N_{i_2}\cdots N_{i_n})^0_0.
\end{array}
\]
这就证明了 $g=0$ 的情形。

其次，对 $g\ge1$，我们将对 $g$ 作归纳来证明定理。对 $g=1$，我们把曲线退化成一个带一个结点的曲线，然后把它正规化。于是由 $g=0$ 的结果有
\[
\begin{array}{rcl}
\dim_{\C}\cV^\dagger_{\lambda(i_1),\dots,\lambda(i_N)}(\mathfrak X)
&=& \displaystyle\sum_{k=1}^{m}\dim_{\C}\cV^\dagger_{\lambda(i_1),\dots,\lambda(i_N),\lambda(k),\lambda(k^\dagger)}(\mathfrak X')\\[10pt]
&=& (N_{i_1}N_{i_2}\cdots N_{i_N}N_kN_{k^\dagger})^0_0.
\end{array}
\]
另一方面我们有
\[
\begin{array}{rcl}
\Bigl(\displaystyle\sum_{k=1}^{m}N_kN_{k^\dagger}\Bigr)^j_i
&=& \displaystyle\sum_{k=1}^{m}\sum_{a=1}^{m}N^a_{k,i}N^j_{a,k^\dagger}
=\sum_{a=1}^{m}\sum_{k=1}^{m}N^a_{i,k}N^k_{a,j^\dagger}\\[10pt]
&=& \displaystyle\sum_{a=1}^{m}\dim\cV_{\lambda(i),\lambda(j^\dagger),\lambda(a),\lambda(a^\dagger)}
(\PP^1;a_1,\dots,a_4;z-a_1,\dots,z-a_4)\\[10pt]
&=& \displaystyle\sum_{a=1}^{m}\sum_{k=1}^{m}N^{k^\dagger}_{i,j^\dagger}N^a_{k^\dagger,a}
=\sum_{k=1}^{m}\Tr(N_{k^\dagger})(N_k)^j_i=W^j_i. \tag{5.35}
\end{array}
\]
因此我们有
\[
\dim_{\C}\cV^\dagger_{\lambda(i_1),\dots,\lambda(i_N)}(\mathfrak X)
=(N_{i_1}N_{i_2}\cdots N_{i_N}W)^0_0.
\]
\hfill$\square$

% [Tutor Note: $g=1$ 那一步的第二个等号漏了求和号 $\sum_{k=1}^{m}$：上一行已对 $k$
%  求和，故第二行应为 $\sum_{k=1}^{m}(N_{i_1}\cdots N_{i_N}N_kN_{k^\dagger})^0_0$
%  $=(N_{i_1}\cdots N_{i_N}\sum_k N_kN_{k^\dagger})^0_0$，这正是 (5.35) 中
%  $\sum_k N_kN_{k^\dagger}=W$ 的来源。照原文排，此处补出正确形式。]

% [原文 p.114]
另外，我们可以改写
\[
\begin{array}{rcl}
\dim_{\C}\cV^\dagger_{i_1,\dots,i_N}(\mathfrak X)
&=&\displaystyle\sum_{k}\sum_{j_1,\dots,j_N}
N^{j_1}_{i_1,k}N^{j_2}_{i_2,j_1}\cdots N^{j_N}_{i_{N-1},j_{N-1}}N^{k}_{i_N,j_N}\\[10pt]
&=&\displaystyle\sum_{k}\Bigl\{\sum_{j_1}N^{j_1}_{k,i_1}
(N_{i_2}\cdots N_{i_{N-1}})^{j_N}_{j_1}N^{k}_{i_N,j_N}\Bigr\}\\[10pt]
&=&\displaystyle\sum_{k}(N_{i_1}N_{i_2}\cdots N_{i_N})^{k}_{k}\\[10pt]
&=&\Tr(N_{i_1}N_{i_2}\cdots N_{i_N}).
\end{array}
\]
这就证明了定理的第二个断言在 $g=1$ 时的情形。

现在设定理对 $g$ 成立，并设 $\mathfrak X$ 是一条带坐标的 $n$ 点亏格 $g+1$ 曲线。令 $\mathfrak X'$ 是一条带坐标的 $n$ 点亏格 $g+1$ 的结点曲线，使得该曲线只有一个结点。令 $\tilde{\mathfrak X}'$ 为由 $\mathfrak X'$ 作正规化得到的带坐标 $n+2$ 点亏格 $g$ 曲线。于是由归纳假设我们有
\[
\begin{array}{rcl}
\dim_{\C}\cV^\dagger_{\lambda(i_1),\dots,\lambda(i_N),\lambda(k),\lambda(k^\dagger)}(\tilde{\mathfrak X}')
&=&(N_{i_1}N_{i_2}\cdots N_{i_N}N_kN_{k^\dagger}W^{g})^{0}_{0}\\[10pt]
&=&\Tr(N_{i_1}N_{i_2}\cdots N_{i_N}N_kN_{k^\dagger}W^{g-1}).
\end{array}
\]
于是由 (5.35) 我们有
\[
\begin{array}{rcl}
\dim_{\C}\cV^\dagger_{i_1,\dots,i_N}(\mathfrak X)
&=&\displaystyle\sum_{k\in P_\ell}
(N_{i_1}N_{i_2}\cdots N_{i_N}N_kN_{k^\dagger}W^{g})^{0}_{0}\\[10pt]
&=&(N_{i_1}N_{i_2}\cdots N_{i_N}N_kN_{k^\dagger}W^{g+1})^{0}_{0}.
\end{array}
\]
这就证明了定理的第一部分。类似地我们有
\[
\sum_{k=1}^{m}\Tr(N_{i_1}N_{i_2}\cdots N_{i_N}N_kN_{k^\dagger}W^{g-1})
=\Tr(N_{i_1}N_{i_2}\cdots N_{i_N}W^{g}).
\]
这就证明了定理的第二部分在 $g+1$ 时的情形。\hfill$\square$

% [Tutor Note: 上面「由 (5.35)」这一步的第二行照原文没有写求和号 $\sum_{k\in P_\ell}$，
%  却仍保留 $N_kN_{k^\dagger}$；正确的读法是 $\sum_{k\in P_\ell}(N_{i_1}\cdots N_{i_N}
%  N_kN_{k^\dagger}W^g)^0_0=(N_{i_1}\cdots N_{i_N}(\sum_k N_kN_{k^\dagger})W^g)^0_0
%  =(N_{i_1}\cdots N_{i_N}W^{g+1})^0_0$，其中最后一个等号正是 (5.35) 给出的
%  $\sum_k N_kN_{k^\dagger}=W$。]

\figomitn{一条 $n$ 点亏格 $g+1$ 曲线 $\mathfrak X$：曲线上画出 $g+1$ 个环柄（下方一排小环）与 $n$ 个标记点，各标记点上方依次标出 $\lambda(i_1),\lambda(i_2),\dots,\lambda(i_k),\dots,\lambda(i_n)$。}

下述定理由 Verlinde 猜测 [V]，并最先由 Moore–Seiberg [MS2] 证明。

% [原文 p.115]
\begin{theorem}[5.22]
$S$-矩阵可同时对角化矩阵 $N_i$，$i=0,\dots,m$。
\end{theorem}

让我们把 $S$-矩阵写成 $S=(S^j_i)$ 的形状，并设
\begin{equation}
S^{-1}N_iS=D_i=(d^{(i)}_j\delta^k_j), \tag{5.36}
\end{equation}
其中矩阵 $N_i$ 的 $(j,k)$-分量是 $N^k_{i,j}$。于是我们有
\[
\sum_l N^l_{i,j}S^k_l=d^{(i)}_kS^k_j.
\]
注意 $N^i_0=1$。于是取 $j=0$ 得到
\[
d^{(i)}_k=\frac{S^k_i}{S^k_0}.
\]
此外，矩阵 $S_i$ 是对称且酉的。（见 (5.29) 与 (5.30)。）因此我们有
\begin{equation}
N^k_{i,j}=\sum_l\frac{S^l_iS^l_jS^l_{k^\dagger}}{S^l_0}. \tag{5.37}
\end{equation}
换言之，由于我们定义了 $N_{ijk}=N^{k^\dagger}_{ij}$，公式 (5.37) 可以写成
\begin{equation}
N_{ijk}=\sum_l\frac{S^l_iS^l_jS^l_k}{S^l_0}. \tag{5.38}
\end{equation}
公式 (5.37) 或 (5.38) 称为 Verlinde 公式。利用 Verlinde 公式，我们可以计算任意亏格的带坐标 $N$ 点曲线的共形块空间的维数。

% [Tutor Note: (5.37) 的来源是把 (5.36) 反解出来：$N_i=S D_iS^{-1}$，故
%  $N^k_{i,j}=\sum_lS^l_j\,d^{(i)}_l\,(S^{-1})^k_l$；由 $S^{-1}=S^\dagger$（酉性，
%  (5.30)）有 $(S^{-1})^k_l=S^{k^\dagger}_l$，再由 $S^j_i=S^i_j$（对称性）写成
%  $S^l_{k^\dagger}$，最后代入 $d^{(i)}_l=S^l_i/S^l_0$ 即得 (5.37)。原文没有写出
%  反解这一步；「对称且酉」二者各用一次，缺一不可。]

\begin{example}[5.23]
对 $\mathfrak g=\mathfrak{sl}(2,\C)$ 与水平 $\ell$，$S$-矩阵由下式给出：
\begin{equation}
S^j_i=\sqrt{\frac{2}{\ell+2}}\,\sin\frac{(i+1)(j+1)\pi}{\ell+2}. \tag{5.39}
\end{equation}
\end{example}
于是由 (5.34)，对亏格 $g$、带坐标的 $N$ 点曲线
\[
\mathfrak X=(C;Q_1,\dots,Q_N;\xi_1,\dots,\xi_N)
\]
我们有
\[
(S^{-1}WS)^j_j=\frac{1}{(S^j_0)^2}.
\]
因此，由定理 5.21 得
\[
\begin{array}{rcl}
\dim_{\C}\cV^\dagger_{\lambda(i_1),\dots,\lambda(i_N)}(\mathfrak X)
&=&\displaystyle\sum_{j=0}^{\ell}(S^j_0)^{2-2g}\prod_{a=1}^{N}\frac{S^j_{i_a}}{S^j_0}\\[10pt]
&=&\displaystyle\Bigl(\frac{\ell+2}{2}\Bigr)^{g-1}
\sum_{j=0}^{\ell}
\frac{\displaystyle\prod_{a=1}^{N}\sin\frac{(i_a+1)(j+1)\pi}{\ell+2}}
{\Bigl(\sin\frac{(j+1)\pi}{\ell+2}\Bigr)^{2g-2+N}}.
\end{array}
\]

% [Tutor Note: $(S^{-1}WS)^j_j=1/(S^j_0)^2$ 这一步：把 (5.34) 的
%  $W=\sum_k(\Tr N_{k^\dagger})N_k$ 共轭对角化，得 $(S^{-1}WS)^j_j
%  =\sum_k(\Tr N_{k^\dagger})\,d^{(k)}_j$，再用 $d^{(k)}_j=S^k_j/S^k_0$
%  与 (5.37) 把 $\Tr N_{k^\dagger}$ 化成 $S^k_{k^\dagger}/S^k_0$ 型求和，
%  用对称性与酉性整理即得。原文把这一整段代数跳过了。]

% [Tutor Note: §5.5 的证明路线，一句话：把因子化定理矩阵化，再把矩阵同时对角化。
%  具体三步：(i) 由 (5.33) 把三点函数看成矩阵 $N_i$ 的元素，于是任意 $N$ 点
%  $\PP^1$ 上共形块的维数就是矩阵乘积的 $(0,0)$ 元（定理 5.21 第一式）；
%  (ii) 引入 $W=\sum_k(\Tr N_{k^\dagger})N_k$ 作为「添一个环柄」的算子，得到
%  亏格 $g$ 的公式 $(N_{i_1}\cdots N_{i_N}W^g)^0_0$（定理 5.21 第二式：对 $g$ 归纳，
%  退化出结点曲线、正规化、再用因子化）；(iii) 由定理 5.22 用 $S$-矩阵同时
%  对角化全部 $N_i$，得到显式求和 (5.37)(5.38)，例 5.23 把 $\mathfrak{sl}(2)$ 的
%  (5.39) 代入即得本节末式。跳步（原文未写出）：定理 5.21 证明中 $W$ 与 $N_i$ 的
%  对易性、以及 $W$ 的幂次在粘合时的移位；定理 5.22 的「同时对角化」只被断言
%  （依据是 Verlinde 代数与融合环的 Grothendieck 环在 $S$ 下同构）；
%  $(S^{-1}WS)^j_j=1/(S^j_0)^2$ 的化简短略。]
